% Exploratory Fig5 artifact: polarity-reversal actuation as a macroscopic
% trapped-charge probe. This is a story test, not publication-final artwork.
\documentclass[border=4pt]{standalone}
\usepackage{polymer-paper-preamble}
\usetikzlibrary{arrows.meta,calc,positioning}

\begin{document}
\resizebox{178mm}{!}{%
\begin{tikzpicture}[
  every node/.style={font=\sffamily\fontsize{6.2}{7.4}\selectfont},
  panelLetter/.style={font=\sffamily\bfseries\fontsize{8}{9.6}\selectfont,
    text=cGray!92!black},
  panelTitle/.style={font=\sffamily\bfseries\fontsize{6.6}{7.9}\selectfont,
    text=cGray!86!black},
  labelStd/.style={font=\sffamily\fontsize{5.5}{6.6}\selectfont,
    text=cGray!88!black},
  labelMute/.style={font=\sffamily\itshape\fontsize{5.1}{6.1}\selectfont,
    text=cGray!62!black},
  wire/.style={cGray!78!black,line width=0.82pt,line cap=round,line join=round},
  apparatus/.style={cGray!18,draw=cGray!78!black,line width=0.82pt},
  polymer/.style={cAmber!10,draw=cAmber!72!black,line width=0.95pt},
  leader/.style={cGray!58!black,line width=0.40pt},
  processArrow/.style={-{Stealth[length=3.8pt,width=2.8pt]},cGray!70!black,
    line width=0.72pt},
  forceToward/.style={-{Stealth[length=4.8pt,width=3.5pt]},cRed!88!black,
    line width=1.02pt},
  forceAway/.style={-{Stealth[length=4.8pt,width=3.5pt]},cRed!88!black,
    line width=1.02pt},
  fieldCue/.style={cGray!48!black,dashed,line width=0.45pt,line cap=round},
  baselineArrow/.style={-{Stealth[length=3.6pt,width=2.6pt]},cGray!62!black,
    line width=0.62pt},
  gapDimension/.style={<->,cGray!55!black,line width=0.48pt},
  qmark/.style={circle,fill=cRed!80!black,draw=cRed!95!black,line width=0.25pt,
    inner sep=0pt,minimum size=1.55mm},
  axis/.style={-{Stealth[length=3.2pt,width=2.4pt]},cGray!72!black,
    line width=0.62pt},
]

% ================================================================
% Panel A — actuation charge, then source-off floating state
% ================================================================
% Panel A
\begin{scope}[shift={(0.00,0.00)}]
  \node[panelLetter,anchor=west] at (0.10,5.12) {a};
  \node[panelTitle,anchor=west] at (0.44,5.12)
    {charge \textrightarrow{} isolate};
  \node[labelMute,anchor=west] at (0.44,4.78)
    {$+5\,\mathrm{kV}$, \textasciitilde20 min};

  % The charge state is the actuation geometry itself: a cantilever faces a
  % nearby drive electrode across an air gap and bends toward it.
  \fill[cGray!34] (0.92,4.05) rectangle (1.82,4.34);
  \draw[apparatus] (0.92,4.05) rectangle (1.82,4.34);
  \draw[wire] (1.37,4.34)--(1.37,4.68);
  \node[labelMute,anchor=east] at (0.82,4.20) {clip};
  \node[labelMute,anchor=west] at (1.52,4.62) {GND};
  % A broad, blunt-ended film strip: the two edges deliberately do not read as
  % a wire or field line at reduced size.
  \fill[polymer,rounded corners=0.35mm]
    (1.18,4.05) .. controls (1.17,3.42) and (1.48,2.86) .. (2.20,2.20)
    -- (2.50,2.50) .. controls (1.86,3.18) and (1.67,3.68) .. (1.64,4.05) -- cycle;
  \draw[cAmber!72!black,line width=1.0pt,rounded corners=0.45mm]
    (1.18,4.05) .. controls (1.17,3.42) and (1.48,2.86) .. (2.20,2.20)
    -- (2.50,2.50) .. controls (1.86,3.18) and (1.67,3.68) .. (1.64,4.05);
  \fill[cGray!28] (3.12,1.78) rectangle (3.42,4.34);
  \draw[apparatus] (3.12,1.78) rectangle (3.42,4.34);
  \node[labelStd,anchor=south,align=center] at (3.27,4.58) {drive electrode};
  \node[labelStd,text=cRed!82!black,anchor=west] at (3.52,4.12)
    {$+5\,\mathrm{kV}$};
  \node[qmark] at (1.38,3.57) {};
  \node[qmark] at (1.55,3.08) {};
  \node[qmark] at (1.88,2.67) {};
  \node[labelStd,text=cRed!80!black,anchor=west] at (0.26,3.04)
    {trapped $q_{\mathrm{tr}}$};
  \draw[leader,cRed!58!black] (0.90,3.08)--(1.45,3.30);

  \draw[gapDimension] (2.45,2.18)--(3.08,2.18);
  \node[labelMute,anchor=south] at (2.76,2.24) {air gap};
  \draw[forceToward] (2.08,3.02)--(2.84,3.02);
  \node[labelStd,text=cRed!84!black,anchor=south] at (2.46,3.08)
    {attraction};

  % The lower rail is an electrical-state transition for the same mounted
  % cantilever, not a second specimen or a motion stage.
  \node[labelStd,anchor=west] at (0.42,1.42) {CHARGE};
  \node[labelMute,anchor=west] at (0.42,1.07) {clip GND};
  \draw[wire] (0.58,0.78)--(1.18,0.78);
  \draw[wire] (0.88,0.78)--(0.88,0.60);
  \draw[wire] (0.67,0.56)--(1.09,0.56);
  \draw[wire] (0.73,0.49)--(1.03,0.49);
  \draw[wire] (0.80,0.42)--(0.96,0.42);
  \draw[processArrow] (1.34,0.86)--(2.26,0.86);
  \node[labelStd,anchor=west] at (2.42,1.42) {OFF / FLOAT};
  \node[labelMute,anchor=west] at (2.42,1.07) {GND open\quad $q_{\mathrm{tr}}$ stays};
  \draw[wire] (2.58,0.78)--(2.88,0.78);
  \draw[wire] (3.02,0.78)--(3.52,0.78);
  \draw[wire] (2.88,0.78)--(3.03,0.95);
  \fill[white] (2.94,0.78) circle (0.52pt);
  \draw[cGray!78!black,line width=0.60pt] (2.94,0.78) circle (0.52pt);
\end{scope}

% ================================================================
% Panel B — force decomposition and polarity switch
% ================================================================
% Panel B
\begin{scope}[shift={(4.45,0.00)}]
  \node[panelLetter,anchor=west] at (0.10,5.12) {b};
  \node[panelTitle,anchor=west] at (0.44,5.12) {polarity changes the force sign};

  % Side-view geometry shared by both polarities.
  \fill[cGray!28] (3.38,1.22) rectangle (3.66,4.26);
  \draw[apparatus] (3.38,1.22) rectangle (3.66,4.26);
  \node[labelStd,anchor=east] at (3.30,2.75) {drive electrode};
  \fill[cGray!38] (0.96,4.12) rectangle (1.98,4.40);
  \draw[apparatus] (0.96,4.12) rectangle (1.98,4.40);
  \draw[wire] (1.47,4.40)--(1.47,4.76);
  \node[labelMute,anchor=west] at (1.72,4.76) {floating clip};
  \fill[polymer,rounded corners=0.35mm]
    (1.24,4.12) .. controls (1.22,3.32) and (0.95,2.52) .. (0.72,1.46)
    -- (1.10,1.38) .. controls (1.54,2.54) and (1.78,3.38) .. (1.70,4.12) -- cycle;
  \draw[cAmber!72!black,line width=1.0pt,rounded corners=0.5mm]
    (1.24,4.12) .. controls (1.22,3.32) and (0.95,2.52) .. (0.72,1.46)
    -- (1.10,1.38) .. controls (1.54,2.54) and (1.78,3.38) .. (1.70,4.12);
  \node[qmark] at (1.35,3.28) {};
  \node[qmark] at (1.17,2.53) {};
  \node[qmark] at (0.97,1.90) {};
  \node[labelStd,text=cRed!80!black,anchor=west] at (0.20,3.52)
    {$q_{\mathrm{tr}}$};
  \draw[leader,cRed!58!black] (0.72,3.46)--(1.18,2.65);

  % At the reversed drive polarity the two force terms are drawn separately
  % from the same specimen position: Maxwell still points toward the driven
  % electrode, while the trapped-charge Coulomb term points away.
  \draw[baselineArrow] (1.74,2.98)--(3.08,2.98);
  \node[labelStd,anchor=south east,text=cGray!74!black] at (3.06,3.04)
    {Maxwell attraction};
  \draw[forceAway] (1.10,2.48)--(0.34,2.48);
  \node[labelStd,text=cRed!84!black,anchor=north west] at (0.02,2.94)
    {Coulomb};
  \node[labelMute,align=center] at (2.22,0.88)
    {$+V$: both terms attract;\\[-0.5pt]after reversal: Coulomb opposes Maxwell};
\end{scope}

% ================================================================
% Panel C — two observed bend directions
% ================================================================
% Panel C
\begin{scope}[shift={(8.90,0.00)}]
  \node[panelLetter,anchor=west] at (0.10,5.12) {c};
  \node[panelTitle,anchor=west] at (0.44,5.12) {same cantilever, opposite bending};

  % The upper observed state bends toward the driven electrode during the
  % grounded +5 kV actuation hold.
  \node[labelStd,anchor=west] at (0.20,4.52) {$+5\,\mathrm{kV}$; clip GND};
  \fill[cGray!28] (3.42,3.34) rectangle (3.66,4.48);
  \draw[apparatus] (3.42,3.34) rectangle (3.66,4.48);
  \fill[cGray!34] (1.08,4.20) rectangle (1.92,4.42);
  \draw[apparatus] (1.08,4.20) rectangle (1.92,4.42);
  \draw[cGray!45!black,dashed,line width=0.62pt]
    (1.50,4.20)--(1.50,3.30);
  \fill[polymer,rounded corners=0.35mm]
    (1.26,4.20) .. controls (1.34,3.80) and (1.70,3.50) .. (2.30,3.18)
    -- (2.58,3.52) .. controls (1.96,3.84) and (1.66,4.10) .. (1.66,4.20) -- cycle;
  \draw[cAmber!72!black,line width=0.92pt]
    (1.26,4.20) .. controls (1.34,3.80) and (1.70,3.50) .. (2.30,3.18)
    -- (2.58,3.52) .. controls (1.96,3.84) and (1.66,4.10) .. (1.66,4.20);
  \node[labelStd,text=cBlue!82!black,anchor=west] at (2.30,3.70) {$+\theta$};

  % After isolation and polarity reversal, the observed cantilever bends away
  % from the same electrode.
  \node[labelStd,anchor=west] at (0.20,2.30) {$-5\,\mathrm{kV}$; clip floating};
  \fill[cGray!28] (3.42,1.12) rectangle (3.66,2.26);
  \draw[apparatus] (3.42,1.12) rectangle (3.66,2.26);
  \fill[cGray!34] (1.08,1.98) rectangle (1.92,2.20);
  \draw[apparatus] (1.08,1.98) rectangle (1.92,2.20);
  \draw[cGray!45!black,dashed,line width=0.62pt]
    (1.50,1.98)--(1.50,1.08);
  \fill[polymer,rounded corners=0.35mm]
    (1.26,1.98) .. controls (1.16,1.58) and (0.92,1.28) .. (0.56,1.02)
    -- (0.86,0.72) .. controls (1.46,1.04) and (1.72,1.48) .. (1.66,1.98) -- cycle;
  \draw[cAmber!72!black,line width=0.92pt]
    (1.26,1.98) .. controls (1.16,1.58) and (0.92,1.28) .. (0.56,1.02)
    -- (0.86,0.72) .. controls (1.46,1.04) and (1.72,1.48) .. (1.66,1.98);
  \node[labelStd,text=cBlue!82!black,anchor=west] at (0.38,1.48) {$-\theta$};
  \node[labelMute,anchor=south east,align=right] at (3.92,0.58)
    {same cantilever; only drive polarity changes};
\end{scope}

% ================================================================
% Panel D — qualitative video-derived trace
% ================================================================
% Panel D
\begin{scope}[shift={(13.35,0.00)}]
  \node[panelLetter,anchor=west] at (0.10,5.12) {d};
  \node[panelTitle,anchor=west] at (0.44,5.12) {video response signature};

  % Axes and zero line.
  \draw[axis] (0.62,1.02)--(0.62,4.08);
  \draw[axis] (0.62,2.52)--(4.02,2.52);
  \node[labelMute,rotate=90,anchor=south] at (0.30,2.56) {bend angle};
  \node[labelMute,anchor=north east] at (4.00,0.92) {time};
  \node[labelMute,anchor=east] at (0.56,3.63) {$+\theta$};
  \node[labelMute,anchor=east] at (0.56,1.46) {$-\theta$};

  % The 20-minute charge hold is preparation, not a calibrated interval on the
  % response axis.  The observed motion starts from its preconditioned state.
  \node[labelMute,anchor=west] at (0.62,4.42)
    {precondition\quad 20 min\quad compressed};
  % The charging hold is outside the observed-response timebase.  The short
  % grey cue and axis break prevent a 20-minute plateau from being compared to
  % the post-reversal mechanical response.
  \draw[cGray!58!black,line width=0.72pt] (0.74,3.62)--(0.98,3.62);
  \draw[cBlue!82!black,line width=1.05pt]
    (1.30,3.62) .. controls (1.52,3.63) and (1.70,3.62) .. (1.88,3.56)
    .. controls (2.06,3.36) and (2.24,2.86) .. (2.42,2.52)
    .. controls (2.60,1.74) and (2.88,1.30) .. (3.16,1.38)
    .. controls (3.58,1.54) and (3.86,2.22) .. (4.02,2.50);
  \draw[white,line width=2.1pt] (1.01,2.52)--(1.23,2.52);
  \draw[cGray!72!black,line width=0.58pt] (1.04,2.42)--(1.12,2.62);
  \draw[cGray!72!black,line width=0.58pt] (1.14,2.42)--(1.22,2.62);
  \draw[cGray!42!black,dashed,line width=0.45pt] (1.88,1.04)--(1.88,4.08);
  \draw[cGray!42!black,dashed,line width=0.55pt] (2.04,1.04)--(2.04,4.08);
  \fill[cRed!85!black] (2.04,3.40) circle (0.75pt);
  \node[labelStd,anchor=south east] at (1.82,4.02) {OFF};
  \node[labelMute,anchor=south east] at (1.82,3.78) {isolate};
  \node[labelStd,text=cRed!84!black,anchor=south west] at (2.10,3.04)
    {reversed drive};
  \draw[gapDimension] (1.88,1.24)--(2.04,1.24);
  \node[labelMute,anchor=north] at (1.96,1.16) {$\sim20\,\mathrm{s}$};
  \node[labelMute,align=center] at (2.05,0.52)
    {continuous bend; slow recovery after reversal};
\end{scope}

% Process separators: guide the eye without turning the figure into cards.
\foreach \x in {4.34,8.79,13.24} {
  \draw[cGray!28!black,dash pattern=on 1.5pt off 2.0pt,line width=0.38pt]
    (\x,0.55)--(\x,4.78);
}
\end{tikzpicture}%
}
\end{document}
